\subsection{Morfismos diferencialmente lisos}
\label{IV.16.10}

\begin{definition}[16.10.1]
\label{IV.16.10.1}
Se dice que un morfismo de preesquemas $f:X\to S$ es
\emph{diferencialmente liso} (o que $X$ es \emph{diferencialmente
liso sobre $S$}) si satisface las condiciones siguientes:
\begin{enumerate}
  \item[(i)] $\Omega_{X/S}^1$ es un $\mathcal{O}_X$-módulo localmente
  proyectivo, es decir, todo punto de $X$ admite una vecindad abierta
  afín $U$ tal que $\Gamma(U,\Omega_{X/S}^1)$ sea un
  $\Gamma(U,\mathcal{O}_X)$-módulo proyectivo (no necesariamente de tipo
  finito).
  \item[(ii)] El homomorfismo canónico (16.3.1.1)
  \[
    \mathbf{S}_{\mathcal{O}_X}^{\bullet}(\Omega_{X/S}^1)
    \longrightarrow\mathscr{G}r_{\bullet}(\mathscr{P}_{X/S})
  \]
  es biyectivo.
\end{enumerate}
En particular, si $\Omega_{X/S}^1$ es localmente libre de rango finito, los
$\mathscr{P}_{X/S}^n$ son $\mathcal{O}_X$-módulos localmente libres de
rango finito (por ser extensiones de tales módulos).
\end{definition}

Se dice que $f$ es \emph{diferencialmente liso en un punto $x\in X$} (o
que $X$ es diferencialmente liso sobre $S$ en el punto $x$) si existe una
vecindad abierta $U$ de $x$ en $X$ tal que $f|U$ sea diferencialmente
liso.

Veremos más adelante (17.12.4) que un morfismo liso es diferencialmente
liso, lo que justifica la terminología; pero el recíproco es falso; en
efecto, un \emph{monomorfismo} $f:X\to S$ es diferencialmente liso,
puesto que $\Omega_{X/S}^1=0$ en virtud de (I, 5.3.8), y por consiguiente
el homomorfismo sobreyectivo (16.3.1.1) es evidentemente biyectivo; ahora
bien, un monomorfismo ni siquiera es necesariamente plano, ni \emph{a fortiori}
liso. Limitémonos aquí a señalar la proposición siguiente:

\begin{proposition}[16.10.2]
\label{IV.16.10.2}
Sean $A$ un anillo, $B$ una $A$-álgebra formalmente lisa para las
topologías discretas (0, 19.3.1). Entonces $\operatorname{Spec}(B)$ es
diferencialmente liso sobre $\operatorname{Spec}(A)$.
\end{proposition}

En efecto, $B\otimes_A B$ es entonces (para las topologías discretas) una
$B$-álgebra formalmente lisa (para uno u otro de los homomorfismos
canónicos $b\mapsto b\otimes1$, $b\mapsto1\otimes b$ de $B$
\oldpage[IV-4]{52}
en $B\otimes_A B$) (0, 19.3.5, (iii)); por tanto, $B\otimes_A B$ es también
una $A$-álgebra formalmente lisa para las topologías discretas (0, 19.3.5,
(ii)). Poniendo $\mathfrak{I}=\mathfrak{I}_{B/A}$, se deduce que
$B\otimes_A B$ es también una $A$-álgebra formalmente lisa para la topología
$\mathfrak{I}$-preádica (0, 19.3.8); como, por hipótesis,
$B=(B\otimes_A B)/\mathfrak{I}$ es una $A$-álgebra formalmente lisa para
las topologías discretas, la proposición resulta de la equivalencia entre a)
y b) en (0, 19.5.4).

\begin{proposition}[16.10.3]
\label{IV.16.10.3}
Para que un morfismo $f:X\to S$ sea diferencialmente liso, es necesario y
suficiente que para todo $x\in X$ exista una vecindad abierta afín $U$ de
$x$, de anillo $A$, tal que $\Gamma(U,\mathscr{P}_{X/S}^{\infty})$ sea una
$A$-álgebra topológica aumentada isomorfa al álgebra completada $\widehat B$,
donde $B=\mathbf{S}_A(V)$, siendo $V$ un $A$-módulo proyectivo y estando
$B$ provista de la topología $B^+$-preádica (donde $B^+$ es el ideal de aumento). Si
$\Omega_{X/S}^1$ es localmente libre de rango finito, se puede sustituir
$\widehat B$ por el álgebra de series formales $A[[T_1,\ldots,T_n]]$.
\end{proposition}

La noción de morfismo diferencialmente liso es evidentemente local sobre
$X$, por lo que podemos reducirnos al caso en que $S=\operatorname{Spec}(B)$,
$X=\operatorname{Spec}(C)$. Consideremos $C\otimes_B C$ como una $C$-álgebra
(para el primer factor); pongamos $\mathfrak{I}=\mathfrak{I}_{C/B}$ y
dotemos $C\otimes_B C$ de la topología $\mathfrak{I}$-preádica; podemos
aplicar a la $C$-álgebra topológica $C\otimes_B C$ y al ideal
$\mathfrak{I}$ de $C\otimes_B C$ la equivalencia entre b) y c) en (0, 19.5.4),
puesto que $(C\otimes_B C)/\mathfrak{I}=C$ es evidentemente una $C$-álgebra
formalmente lisa para las topologías discretas. La topología sobre
$\Gamma(U,\mathscr{P}_{X/S}^{\infty})$ es evidentemente la topología de límite
proyectivo de este anillo (16.1.11).

Se observará que el entero $n$ del enunciado de (16.10.3) es el \emph{rango} de
$\Omega_{X/S}^1$ en el punto $x$. Veremos más adelante (17.13.5) que cuando
$f$ es diferencialmente liso y localmente de tipo finito, $n$ es también
igual a la dimensión de la fibra $f^{-1}(f(x))$ en el punto $x$.

\begin{proposition}[16.10.4]
\label{IV.16.10.4}
Sean $f:X\to S$, $g:S'\to S$ dos morfismos, y pongamos
$X'=X\times_S S'$, $f'=f_{(S')}:X'\to S'$.
\begin{enumerate}
  \item[(i)] Si $f$ es diferencialmente liso, lo mismo ocurre con $f'$.
  \item[(ii)] Recíprocamente, si $g$ es fielmente plano y cuasicompacto, y si
  $f'$ es diferencialmente liso y $\Omega_{X'/S'}^1$ es un
  $\mathcal{O}_{X'}$-módulo de tipo finito, $f$ es diferencialmente liso y
  $\Omega_{X/S}^1$ es un $\mathcal{O}_X$-módulo de tipo finito.
\end{enumerate}
\end{proposition}

En efecto, si $f$ es diferencialmente liso, los
$\mathscr{G}r_n(\mathscr{P}_{X/S})$ son $\mathcal{O}_X$-módulos
\emph{planos}; por consiguiente, en virtud de (16.4.6), el homomorfismo
\[
  \mathscr{G}r_n(\mathscr{P}_{X/S})\otimes_{\mathcal{O}_X}\mathcal{O}_{X'}
  \longrightarrow\mathscr{G}r_n(\mathscr{P}_{X'/S'})
\]
es biyectivo para todo $n$, y en virtud de la conmutatividad del diagrama
(16.2.1.3), se deduce de la definición (16.10.1) que $f'$ es
diferencialmente liso. Por otra parte, si $g$ es fielmente plano y
cuasicompacto, se deduce también de (16.4.6) que
\[
  \mathscr{G}r_n(\mathscr{P}_{X/S})\otimes_{\mathcal{O}_X}\mathcal{O}_{X'}
  \longrightarrow\mathscr{G}r_n(\mathscr{P}_{X'/S'})
\]
es biyectivo para todo $n$. Supongamos entonces que $f'$ es diferencialmente liso y
$\Omega_{X'/S'}^1$ tiene rango finito. Como la proyección canónica $X'\to X$ es
un morfismo fielmente plano y cuasicompacto, se deduce primero de (2.5.2)
que $\Omega_{X/S}^1$ es un $\mathcal{O}_X$-módulo localmente libre de rango
finito, y luego de (2.2.7) que el homomorfismo canónico (16.3.1.1) es biyectivo,
por lo que $f$ es diferencialmente liso.

\begin{proposition}[16.10.5]
\label{IV.16.10.5}
Para que un morfismo localmente de tipo finito $f:X\to S$ sea
diferencialmente liso, es necesario y suficiente que la inmersión diagonal
$\Delta_f:X\to X\times_S X$ sea cuasirregular.
\end{proposition}

\oldpage[IV-4]{53}

La cuestión es local, por lo que podemos reducirnos al caso en que $S$ y $X$
son afines, y por consiguiente el subpreesquema diagonal de $X\times_S X$ es
cerrado. La hipótesis de que $f$ es localmente de tipo finito implica que
$\Delta_f$ es localmente de presentación finita (I, 4.3.1), por lo que el
subpreesquema diagonal de $X\times_S X$ está definido por un ideal de tipo
finito $\mathscr{I}$, y $\Omega_{X/S}^1=\mathscr{I}/\mathscr{I}^2$ es un
$\mathcal{O}_X$-módulo de tipo finito. La proposición resulta entonces
inmediatamente de la comparación de las condiciones en (16.10.1) y (16.9.4).

\begin{remark}[16.10.6]
\label{IV.16.10.6}
Sea $f:X\to S$ un morfismo tal que el $\mathcal{O}_X$-módulo
$\Omega_{X/S}^1$ sea localmente libre de rango finito. Se deduce entonces de
(0, 20.4.7) que todo $x\in X$ posee una vecindad abierta $U$ tal que existe
una familia finita $(z_\lambda)_{\lambda\in L}$ de secciones de
$\mathcal{O}_X$ sobre $U$ para las cuales $(dz_\lambda)_{\lambda\in L}$
forma una \emph{base} del $\Gamma(U,\mathcal{O}_X)$-módulo
$\Gamma(U,\Omega_{X/S}^1)$.
\end{remark}
